% ID: MAT-TRI-EQG-002
% Titolo: Equazioni goniometriche e periodicità (3)
% Disciplina: Matematica
% Area: trigonometria
% Argomento: equazioni_goniometriche
% Classe: Quarta secondaria di secondo grado
% Difficolta: 4
% Tipo: problema_numerico
% Tag: matematica, trigonometria, equazioni_goniometriche
% Fonte: verifica_goniometria_avanzata_anonimizzata
% Provenienza: materiale fornito dal docente; testo anonimizzato; soluzione redatta per Exergo

\begin{esercizio}
Risolvi le seguenti equazioni: 

\begin{enumerate}
\item $\displaystyle
\sin\left(2x+\dfrac{\pi}{3}\right)
=\cos\left(x-\dfrac{\pi}{3}\right);$

\item $\displaystyle
\sqrt{3}\,\sin 2x=2\cos^2x.$
\end{enumerate}
Salvo diversa indicazione, lavora in $\mathbb R$; angoli in radianti.
\end{esercizio}

\begin{soluzione}
a) Scrivendo il secondo membro come $\sin(5\pi/6-x)$, si ottiene
$2x+\pi/3=5\pi/6-x+2k\pi$ oppure $2x+\pi/3=\pi/6+x+2k\pi$.
Quindi $x=\pi/6+2k\pi/3$ oppure $x=-\pi/6+2k\pi$.
b) $2\cos x(\sqrt3\sin x-\cos x)=0$:
$x=\pi/2+k\pi$ oppure $x=\pi/6+k\pi$. In tutte le famiglie di soluzioni $k\in\mathbb{Z}$.
\end{soluzione}
